\section{The induced congruence topology}\label{sec:completion}

For $m\ge1$, define a congruence on $F(D)$ by
\[
 \theta_m=
 \bigcap\{\ker(\pi_T):
 T\text{ is a finite-complement compatible extension and }
 |T\setminus A|\le m\}.
\]
Equivalently, $x\equiv y\pmod{\theta_m}$ when every such extension identifies
$x$ and $y$.  Each $\theta_m$ is a congruence and
\[
 \theta_{m+1}\subseteq\theta_m.
\]
By \cref{cor:finite-complement-separation},
\[
 \bigcap_{m\ge1}\theta_m=\Delta_{F(D)}.
\]
Hence the cosets of the $\theta_m$ form a countable Hausdorff congruence
topology on $F(D)$; this is a standard construction for universal algebras
\cite{BulmanFleming1971}.  We denote its Hausdorff completion by
$\widehat{F(D)}_{\mathrm{fc}}$.

\subsection{Finite bases}

When $A$ is finite, the preceding topology is not a new kind of completion.

\begin{proposition}[Finite base gives the profinite topology]
\label{prop:finite-base-profinite}
Assume that $A$ and $\Sigma$ are finite.  Then the topology generated by the
congruences $\theta_m$ is the ordinary profinite topology on $F(D)$.
\end{proposition}

\begin{proof}
First, each $\theta_m$ has finite index.  Up to relabelling of the finitely many
points outside the fixed finite set $A$, there are only finitely many total
$\Sigma$-algebra structures on carriers of size at most $|A|+m$.  Thus only
finitely many kernels occur in the defining intersection for $\theta_m$, and
an intersection of finitely many finite-index congruences again has finite
index.  Hence every $\theta_m$ is an entourage for the profinite topology.

Conversely, let $\alpha$ be any finite-index congruence on $F(D)$.  Apply
\cref{lem:finite-complement-rees} to the finite set $X=A$.  This gives a
finite-index congruence $\beta$ whose restriction to $A$ is equality.  Put
\[
 \gamma=\alpha\cap\beta.
\]
Then $\gamma$ has finite index, $\gamma\subseteq\alpha$, and $A$ remains
pointwise distinct in $F(D)/\gamma$.  After identifying the image of $A$ with
$A$ itself, the quotient $F(D)/\gamma$ is a finite-complement compatible
extension.  If its complement has size $m$, then
\[
 \theta_m\subseteq\gamma\subseteq\alpha.
\]
Thus the family $(\theta_m)$ is cofinal among the finite-index congruences, so
it induces exactly the profinite topology.
\end{proof}

The remainder of the section concerns the situation in which $A$ may be
infinite.  Then an allowed target $T$ can be infinite, and the congruences
$\theta_m$ need not have finite index.

\subsection{A criterion over an arbitrary base}

We use one elementary observation about finite deterministic orbits.

\begin{lemma}\label{lem:finite-orbit-factorials}
Let $s_0,s_1,\ldots$ be an orbit of a self-map.  If the orbit takes at most
$N$ distinct values, then $s_{k!}$ is independent of $k$ for all sufficiently
large $k$.  The threshold can be chosen using only $N$.
\end{lemma}

\begin{proof}
Among the first $N+1$ positions two coincide.  The orbit is therefore eventually
periodic, with preperiod and period bounded by $N$; every sufficiently large
factorial lies beyond the preperiod and is divisible by the period.
\end{proof}

A \emph{one-hole unary context} $C[-]$ is a term with one distinguished
occurrence of a variable.  We say that it \emph{fixes the base} if substitution
of any $a\in A$ is defined in $D$ and evaluates to $a$:
\[
 C^D(a)=a\qquad(a\in A).
\]
We call the context nontrivial when it contains at least one operation symbol.

\begin{theorem}[Infinite-base completion criterion]
\label{thm:infinite-base-criterion}
Suppose that some basic application in $D$ is undefined, and let $c_0\in
F(D)\setminus A$ be the corresponding normal term.  If there is a nontrivial
one-hole unary context $C[-]$ that fixes the base, then the canonical map
\[
 F(D)\longrightarrow\widehat{F(D)}_{\mathrm{fc}}
\]
is not surjective.
\end{theorem}

\begin{proof}
Define
\[
 c_{n+1}=C[c_n]\qquad(n\ge0).
\]
Because $C[-]$ is nontrivial and $c_0$ is non-base, no normalization step can
remove the occurrence of $c_n$ inside $C[c_n]$; hence the construction depth
strictly increases and the terms $c_n$ are pairwise distinct.

Consider the subsequence $d_k=c_{k!}$.  Fix $m$ and a finite-complement
compatible extension $T$ with $|T\setminus A|\le m$.  The images of the $c_n$
form an orbit under the unary operation induced by $C[-]$ on $T$.  If this
orbit enters $A$, it is constant from then on because $C$ fixes every element
of $A$.  Otherwise it remains in the at most $m$ elements of $T\setminus A$.
Thus in either case the orbit has at most $m+1$ values before becoming
periodic.  By \cref{lem:finite-orbit-factorials}, for all sufficiently large
$k,l$, with a threshold depending only on $m$,
\[
 \pi_T(d_k)=\pi_T(d_l).
\]
Since the threshold is uniform over all such $T$, we have
$d_k\equiv d_l\pmod{\theta_m}$ for all sufficiently large $k,l$.  As $m$ was
arbitrary, $(d_k)$ is Cauchy.

It remains to show that this Cauchy sequence has no limit in $F(D)$.  Fix
$x\in F(D)$.  Choose $N$ so large that the construction depth of $c_N$ exceeds
that of $x$, and apply \cref{lem:finite-complement-rees} to
$X=\{x,c_N\}$.  In the resulting quotient, $x$ and every selected subterm are
singleton classes.  The next term $c_{N+1}=C[c_N]$ is not among those selected
subterms: it has greater depth than both selected roots.  Hence it belongs to
the collapsed ideal class.  Since that class is absorbing under every basic
operation, all later $c_j$ also belong to it.  Therefore this one
finite-complement quotient separates $x$ from every sufficiently late $d_k$.
So $(d_k)$ cannot converge to $x$.

Since $x$ was arbitrary, the Cauchy sequence $(d_k)$ has no limit in $F(D)$.
Its limit in the Hausdorff completion is therefore a point outside the image of
$F(D)$.
\end{proof}

\begin{remark}
The theorem is only a sufficient criterion.  It does not characterize when the
finite-complement congruence topology over an infinite base is incomplete.
The point of the hypothesis $C^D|_A=\mathrm{id}_A$ is that, although the target
algebra may be infinite, the particular orbit used in the proof cannot employ
the retained base as unbounded memory.
\end{remark}
